% Partie : logique du second ordre
% Chapitre : syntaxe et sémantique
% Section : langage de la logique du second ordre

\documentclass[../../../include/open-logic-section]{subfiles}

\begin{document}

\olfileid[fr]{sol}{syn}{lan}

\olsection{Le langage de la logique du second ordre}

Comme en logique du premier ordre, les expressions
se forment à partir d'un vocabulaire de base
comprenant des \emph{!!{variable}s}, des
\emph{!!{constant}s}, des \emph{!!{predicate}s}
et parfois des \emph{!!{function}s}. Avec les
connecteurs logiques, les quantificateurs et les
signes de ponctuation, notamment les parenthèses
et les virgules, on construit des \emph{termes}
et des \emph{!!{formula}s}. La nouveauté est que
la logique du second ordre possède aussi des
variables de relation et de fonction, sur lesquelles
on peut quantifier. À ses symboles logiques du
premier ordre s'ajoutent donc :

\begin{enumerate}
\item Pour chaque arité~$n$ admise, un ensemble
  infini dénombrable de !!{variable}s de relation
  du second ordre :
  $\Obj V_0^n$, $\Obj V_1^n$, $\Obj V_2^n$, \dots
\item Pour chaque arité~$n$ admise, un ensemble
  infini dénombrable de !!{variable}s de fonction
  du second ordre :
  $\Obj u_0^n$, $\Obj u_1^n$, $\Obj u_2^n$, \dots
\end{enumerate}

En logique du premier ordre, on inclut
habituellement l'!!{identity}~$\eq$.
Les symboles non logiques d'un langage~$\Lang{L}$
permettent de désigner des relations, fonctions
et objets particuliers. Certaines !!{sentence}s
n'en contiennent aucun, mais le seul symbole
d'égalité exprime peu de propriétés spécifiques.
En logique du second ordre, les variables de
relation et de fonction fournissent, à chaque
arité, des places pour ces symboles : on peut
ainsi reformuler des schémas d'énoncés du premier
ordre sans fixer d'avance de vocabulaire non
logique. Pour retrouver l'interprétation d'un
symbole dans une structure donnée, on affecte à
la variable libre correspondante la relation ou
la fonction que ce symbole désigne. Quantifier
sur cette variable fait varier les relations ou
les fonctions considérées et change en général
les conditions de vérité.

\end{document}
